% !PRINTED 700
\weberpage{700}
%<<CONT
1. \gesperrt{The polygon quotient (1) is a differential,
if the classes $A$, $B$ of $\mathfrak{A}$ and $\mathfrak{B}$ are relatively prime
and satisfy the condition:}
\weq{2}{A = WB.}

If $a$ and $b$ are the order numbers of $A$ and $B$, then:
\weq{3}{a - b = 2p - 2,}
and since $A$ is in any case of the second kind, by \S\,199, (15):
\weq{4}{(O, A) = b + p - 1.}

If $B$ consists of a single point $\mathfrak{P}$, then $(O, A) = p$,
and the dimension of $A$ is just as great as that of $W$; if therefore
\[
(\mathfrak{W}_1, \mathfrak{W}_2, \ldots, \mathfrak{W}_p)
\]
is a basis of $W$, then
\[
(\mathfrak{P} \mathfrak{W}_1, \mathfrak{P} \mathfrak{W}_2, \ldots, \mathfrak{P} \mathfrak{W}_p)
\]
is a basis of $A$. Hence $A$ is an improper class with
the divisor $\mathfrak{P}$, and it follows:

2. \gesperrt{A single point $\mathfrak{P}$ cannot be the lower
polygon of a differential.}

If the class $A$ has a divisor $\mathfrak{M}$ of degree $m$, then $\mathfrak{M}$ must
be contained in the divisor of $B$; for $\mathfrak{M}$ must divide each of the
polygons $\mathfrak{W}\mathfrak{B}$, and consequently, since $W$ is a proper class,
each of the $\mathfrak{B}$. If
\[
\mathfrak{M} A' = A, \quad \mathfrak{M} B' = B,
\]
then
\weq{5}{(O, A) = (O, A')}
must hold. But by (3) and (4), applied to $A$ and $A'$,
if $A'$ is of the second kind:
\[
(O, A) = (O, A') + m,
\]
hence $(O, A) > (O, A')$; consequently $m = 0$ must hold.

But $A'$ can be of the first kind only if
$B$ is an isolated polygon and $B = \mathfrak{M}$. Then:
\[
A' = W, \quad (O, A') = p,
\]
so that $(O, A) = p$, $b = 1$ must hold, and we come to the
case of theorem 2. Consequently:

3. \gesperrt{Every polygon of more than one point can
be the lower polygon of a differential.}

By a \gesperrt{differential of the second kind} one understands
one whose lower polygon is a power of a single point $\mathfrak{P}$
%CONT>>
